\documentclass[12pt,a4paper]{article}
\usepackage[margin=2.5cm]{geometry}
\usepackage{amsmath,amssymb,amsfonts}
\usepackage{physics}
\usepackage{enumitem}
\usepackage{fancyhdr}
\usepackage{lastpage}
\usepackage{graphicx}
\usepackage{multicol}
\usepackage{array}
\usepackage{tabularx}
\usepackage{tikz}

% Page style
\pagestyle{fancy}
\fancyhf{}
\renewcommand{\headrulewidth}{0.4pt}
\renewcommand{\footrulewidth}{0.4pt}
\lhead{PHYS10012}
\rhead{Mock Examination}
\cfoot{Page \thepage\ of \pageref{LastPage}}

% Question numbering
\newcounter{questioncounter}
\newcommand{\question}[1]{%
  \stepcounter{questioncounter}%
  \vspace{0.5cm}%
  \noindent\textbf{Question \thequestioncounter.} \hfill [#1 marks]%
  \vspace{0.2cm}%
  \par%
}

% Sub-part environments
\newlist{parts}{enumerate}{2}
\setlist[parts,1]{label=(\alph*),leftmargin=2em,itemsep=0.3cm}
\setlist[parts,2]{label=(\roman*),leftmargin=2em,itemsep=0.2cm}

% Mark allocation in sub-parts
\renewcommand{\marks}[1]{\hfill [#1]}

% MCQ environment
\newcounter{mcqcounter}
\newcommand{\mcq}[1]{%
  \stepcounter{mcqcounter}%
  \vspace{0.3cm}%
  \noindent\textbf{\themcqcounter.} #1 \hfill [2 marks]%
  \vspace{0.1cm}%
}
\newlist{choices}{enumerate}{1}
\setlist[choices]{label=\Alph*.,leftmargin=2.5em,itemsep=0.1cm}

% Dirac notation shortcuts
\renewcommand{\ket}[1]{\left|#1\right\rangle}
\renewcommand{\bra}[1]{\left\langle #1\right|}
\renewcommand{\braket}[2]{\left\langle #1\middle|#2\right\rangle}
\renewcommand{\ketbra}[2]{\left|#1\right\rangle\!\left\langle #2\right|}
\renewcommand{\expect}[1]{\left\langle #1\right\rangle}

\begin{document}

% Title block
\begin{center}
{\Large\textbf{UNIVERSITY OF BRISTOL}}\\[0.3cm]
{\large Faculty of Science}\\[0.5cm]
{\Large\textbf{MOCK EXAMINATION --- Paper 8 (Extreme)}}\\[0.3cm]
{\large\textbf{Core Physics I --- Classical, Quantum and Thermal Physics}}\\[0.2cm]
{\large Module Code: PHYS10012}\\[0.5cm]
{\large Duration: 3 Hours}\\[0.3cm]
{\large Total Marks: 100}\\[0.5cm]
\end{center}

\noindent\rule{\textwidth}{0.4pt}

\vspace{0.3cm}
\noindent\textbf{Instructions:}
\begin{itemize}[itemsep=0.1cm]
\item Answer \textbf{ALL} questions.
\item \textbf{Section A} (Oscillations and Waves): 10 multiple-choice questions, 20 marks total.
\item \textbf{Section B} (Quantum Physics): 10 multiple-choice questions, 20 marks total.
\item \textbf{Section C}: 4 long-answer questions, 60 marks total.
\item An approved calculator is permitted.
\item A formula sheet is provided at the end of this paper.
\item Show all working for Section C. Marks will be awarded for clear, logical solutions.
\end{itemize}

\noindent\rule{\textwidth}{0.4pt}
\newpage

% ============================================================
% SECTION A: OSCILLATIONS AND WAVES MCQ
% ============================================================
\section*{Section A: Oscillations and Waves}
\noindent\textit{Answer all 10 questions. Each question is worth 2 marks.}\\[0.3cm]

\setcounter{mcqcounter}{0}

% QUESTION A1: Non-linear SHM from cosine potential
\mcq{A particle of mass $m$ moves in the potential $V(x) = V_0\!\left(1 - \cos\!\left(\dfrac{x}{a}\right)\right)$, where $V_0$ and $a$ are positive constants. For small displacements about $x = 0$, what is the angular frequency $\omega$ of oscillation?}
\begin{choices}
\item $\displaystyle\sqrt{\frac{V_0}{ma^2}}$
\item $\displaystyle\sqrt{\frac{2V_0}{ma^2}}$
\item $\displaystyle\sqrt{\frac{V_0}{2ma^2}}$
\item $\displaystyle\sqrt{\frac{V_0}{ma}}$
\end{choices}

% QUESTION A2: Damped oscillation energy
\mcq{A damped harmonic oscillator has damping constant $\gamma$ and initial energy $E_0$. The amplitude decays as $A(t) = A_0\,e^{-\gamma t/2}$. Which expression gives the energy $E(t)$ at time $t$?}
\begin{choices}
\item $E_0\,e^{-\gamma t/2}$
\item $E_0\,e^{-\gamma t}$
\item $E_0\,e^{-2\gamma t}$
\item $E_0\,(1 - e^{-\gamma t})$
\end{choices}

% QUESTION A3: Forced oscillation FWHM
\mcq{For a lightly damped forced oscillator with damping constant $\gamma$, the full-width at half-maximum (FWHM) of the \emph{power absorption} curve (plotted against driving frequency $\omega$) is approximately:}
\begin{choices}
\item $\gamma/2$
\item $\gamma$
\item $2\gamma$
\item $\omega_0/Q$
\end{choices}

% QUESTION A4: Counter-propagating waves of different amplitudes
\mcq{Two harmonic waves of the same frequency and wavenumber travel in opposite directions on a string: $y_1 = A_1\sin(kx - \omega t)$ and $y_2 = A_2\sin(kx + \omega t)$, where $A_1 \neq A_2$. The superposition $y_1 + y_2$ is:}
\begin{choices}
\item A pure standing wave.
\item A pure travelling wave.
\item A standing wave superimposed with a travelling wave.
\item Zero everywhere.
\end{choices}

% QUESTION A5: Quarter-wave impedance transformer
\mcq{Three strings under the same tension are joined in series: string~1 (impedance $Z_1$), a middle section of length $\ell$ (impedance $Z_m$), and string~3 (impedance $Z_3$). If the middle section has length $\ell = \lambda/4$ (a quarter-wavelength at the operating frequency), what value of $Z_m$ eliminates all reflection from an incident wave in string~1?}
\begin{choices}
\item $Z_m = (Z_1 + Z_3)/2$
\item $Z_m = Z_1 Z_3 / (Z_1 + Z_3)$
\item $Z_m = \sqrt{Z_1 Z_3}$
\item $Z_m = Z_3$
\end{choices}

% QUESTION A6: Coupled oscillators matrix
\mcq{Two identical masses $m$ are connected by springs of constant $k$ to fixed walls, and by a coupling spring $k_c$ to each other. In the matrix equation $m\ddot{\mathbf{x}} = -\mathbf{K}\mathbf{x}$ (where $\mathbf{x} = (x_1,\,x_2)^T$), which matrix $\mathbf{K}$ gives the correct equations of motion?}
\begin{choices}
\item $\displaystyle\begin{pmatrix} k + k_c & -k_c \\ -k_c & k + k_c \end{pmatrix}$
\item $\displaystyle\begin{pmatrix} k & k_c \\ k_c & k \end{pmatrix}$
\item $\displaystyle\begin{pmatrix} k + k_c & k_c \\ k_c & k + k_c \end{pmatrix}$
\item $\displaystyle\begin{pmatrix} k & -k_c \\ -k_c & k \end{pmatrix}$
\end{choices}

% QUESTION A7: Non-linear dispersion — wave packet spreading
\mcq{A wave packet propagates in a medium with dispersion relation $\omega(k)$. The wave packet will spread (disperse) over time if:}
\begin{choices}
\item $\dfrac{d\omega}{dk} \neq 0$
\item $\dfrac{d^2\omega}{dk^2} \neq 0$
\item $\dfrac{d\omega}{dk} = 0$
\item $\omega/k$ is constant for all $k$
\end{choices}

% QUESTION A8: Relativistic Doppler vs classical
\mcq{A source emitting frequency $f_0$ moves toward a stationary observer at speed $v_s \ll c$. The classical Doppler formula gives $f'_\text{cl} = f_0/(1 - v_s/c)$ and the relativistic formula gives $f'_\text{rel} = f_0\sqrt{(1+\beta)/(1-\beta)}$ where $\beta = v_s/c$. To first order in $\beta$, the relativistic correction modifies the classical result by a factor of approximately:}
\begin{choices}
\item $1 - \beta/2$
\item $1 + \beta/2$
\item $1 + \beta$
\item There is no first-order correction; they agree to first order in $\beta$.
\end{choices}

% QUESTION A9: Energy in standing wave vs travelling wave
\mcq{A standing wave $y = 2A\sin(kx)\cos(\omega t)$ is formed on a string of length $L = n\lambda/2$ (i.e.\ $n$ half-wavelengths fit). How does its total energy compare to that of a single travelling wave $y = A\sin(kx - \omega t)$ on the same string?}
\begin{choices}
\item The standing wave has twice the total energy.
\item The standing wave has the same total energy.
\item The standing wave has half the total energy.
\item The energies cannot be compared without knowing $n$.
\end{choices}

% QUESTION A10: Mach cone angle
\mcq{An object moves through air at Mach~2 (i.e.\ at twice the speed of sound). The half-angle $\theta$ of the Mach cone (shock wave cone) is:}
\begin{choices}
\item $60^\circ$
\item $45^\circ$
\item $30^\circ$
\item $15^\circ$
\end{choices}

\newpage

% ============================================================
% SECTION B: QUANTUM PHYSICS MCQ
% ============================================================
\section*{Section B: Quantum Physics}
\noindent\textit{Answer all 10 questions. Each question is worth 2 marks.}\\[0.3cm]

\setcounter{mcqcounter}{0}

% QUESTION B1: Commutator algebra
\mcq{If $[\hat{A},\hat{B}] = i\hat{C}$ for operators $\hat{A}$, $\hat{B}$, and $\hat{C}$, what is $[\hat{B},\hat{A}]$?}
\begin{choices}
\item $i\hat{C}$
\item $-i\hat{C}$
\item $\hat{C}$
\item $0$
\end{choices}

% QUESTION B2: Compatible observables
\mcq{If two Hermitian operators $\hat{A}$ and $\hat{B}$ satisfy $[\hat{A},\hat{B}] = 0$, which statement is correct?}
\begin{choices}
\item Measuring $\hat{A}$ always gives the same result as measuring $\hat{B}$.
\item $\hat{A}$ and $\hat{B}$ share a common set of eigenstates, and both can be measured simultaneously without mutual disturbance.
\item $\hat{A}$ and $\hat{B}$ must be equal.
\item The uncertainty product $\Delta A\,\Delta B$ must equal $\hbar/2$.
\end{choices}

% QUESTION B3: Unitary diagonalisation
\mcq{Let $U$ be the unitary matrix whose columns are the normalised eigenvectors of a Hamiltonian $H$. Then $U^\dagger H\, U$ equals:}
\begin{choices}
\item The identity matrix $I$.
\item $H$ itself.
\item A diagonal matrix with the eigenvalues of $H$ on the diagonal.
\item The zero matrix.
\end{choices}

% QUESTION B4: No-cloning theorem
\mcq{The no-cloning theorem states that it is impossible to create an exact copy of an arbitrary unknown quantum state. This is fundamentally because:}
\begin{choices}
\item Copying requires infinite energy.
\item The Heisenberg uncertainty principle forbids knowing the state precisely.
\item Quantum evolution is linear (unitary), and linearity is incompatible with a universal cloning operation.
\item Entanglement prevents the state from being localised.
\end{choices}

% QUESTION B5: Density matrix of pure state
\mcq{For a system in a pure state $\ket{\psi}$, the density matrix is $\rho = \ketbra{\psi}{\psi}$. What is $\text{Tr}(\rho^2)$?}
\begin{choices}
\item $0$
\item $1/2$
\item $1$
\item It depends on the state $\ket{\psi}$.
\end{choices}

% QUESTION B6: Time-energy uncertainty
\mcq{Two energy eigenstates of a system have energies $E_1$ and $E_2$. A superposition of these states will exhibit oscillatory behaviour. The characteristic timescale $\tau$ of these oscillations is:}
\begin{choices}
\item $\tau = \hbar / (E_1 + E_2)$
\item $\tau = \hbar / |E_1 - E_2|$
\item $\tau = h / (E_1 E_2)$
\item $\tau = (E_1 - E_2) / \hbar$
\end{choices}

% QUESTION B7: Singlet state correlations
\mcq{Two spin-$\frac{1}{2}$ particles are in the singlet state $\ket{\Psi^-} = \frac{1}{\sqrt{2}}\bigl(\ket{\!\uparrow\downarrow} - \ket{\!\downarrow\uparrow}\bigr)$. What is the expectation value $\expect{S_{1z}\,S_{2z}}$?}
\begin{choices}
\item $+\hbar^2/4$
\item $-\hbar^2/4$
\item $0$
\item $-\hbar^2/2$
\end{choices}

% QUESTION B8: Partial trace of Bell state
\mcq{Two qubits are in the Bell state $\ket{\Phi^+} = \frac{1}{\sqrt{2}}\bigl(\ket{\!\uparrow\uparrow} + \ket{\!\downarrow\downarrow}\bigr)$. The reduced density matrix $\rho_1$ of qubit~1, obtained by tracing over qubit~2, is:}
\begin{choices}
\item $\ketbra{\!\uparrow}{\uparrow\!}$
\item $\frac{1}{2}\ketbra{\!\uparrow}{\uparrow\!} + \frac{1}{2}\ketbra{\!\downarrow}{\downarrow\!}$
\item $\frac{1}{2}\bigl(\ketbra{\!\uparrow}{\uparrow\!} + \ketbra{\!\uparrow}{\downarrow\!} + \ketbra{\!\downarrow}{\uparrow\!} + \ketbra{\!\downarrow}{\downarrow\!}\bigr)$
\item $\frac{1}{\sqrt{2}}\bigl(\ket{\!\uparrow} + \ket{\!\downarrow}\bigr)\frac{1}{\sqrt{2}}\bigl(\bra{\!\uparrow} + \bra{\!\downarrow}\bigr)$
\end{choices}

% QUESTION B9: Fermion ground state energy
\mcq{$N$ non-interacting spin-$\frac{1}{2}$ fermions are placed in a one-dimensional infinite square well of width $L$. Each single-particle energy level $E_n = n^2 E_1$ (where $E_1 = \pi^2\hbar^2/(2mL^2)$) can hold at most 2 fermions (one spin-up, one spin-down). The ground-state energy of the system for $N = 6$ fermions is:}
\begin{choices}
\item $6E_1$
\item $14E_1$
\item $28E_1$
\item $42E_1$
\end{choices}

% QUESTION B10: Parity violation
\mcq{Which discrete symmetry is violated in weak interactions but preserved in electromagnetic and strong interactions?}
\begin{choices}
\item Charge conjugation ($C$) only.
\item Parity ($P$).
\item Time reversal ($T$) only.
\item The combined symmetry $CPT$.
\end{choices}

\newpage

% ============================================================
% SECTION C: LONG ANSWER QUESTIONS
% ============================================================
\section*{Section C: Long Answer Questions}
\noindent\textit{Answer all 4 questions. Show all working.}

\setcounter{questioncounter}{0}

% ---- QUESTION C1: MECHANICS (15 marks) ----
\question{15}

\noindent\textbf{Snooker Ball Impulse, Rolling Friction, and Loop-the-Loop}

\begin{parts}
\item A uniform solid sphere (billiard ball) of mass $M$ and radius $R$ rests on a horizontal surface with coefficient of kinetic friction $\mu_k$. A cue strikes the ball with a horizontal impulse $J$ at a height $h$ above the centre of the ball.
\begin{parts}
\item Find the initial translational velocity $v_0$ and the initial angular velocity $\omega_0$ of the ball immediately after the impulse. \qmarks{2}
\item Show that if the cue strikes at height $h = \frac{2}{5}R$ above the centre, the ball immediately rolls without slipping (i.e.\ $v_0 = R\omega_0$), so no friction is needed. At what height above the \emph{table} (i.e.\ above the lowest point of the ball) should the cue strike? \qmarks{3}
\end{parts}

\item For a general strike height $h \neq \frac{2}{5}R$, the ball initially slides. Friction acts until rolling is achieved.
\begin{parts}
\item Write the equations of motion for the translational and rotational motion during the sliding phase. \qmarks{2}
\item Find the time $t_\text{roll}$ at which rolling without slipping begins and the velocity $v_\text{roll}$ at that instant. For a centre-strike ($h = 0$), show that $v_\text{roll} = \frac{5}{7}\,v_0$. \qmarks{3}
\end{parts}

\item A solid disk of mass $M = 0.5\;\text{kg}$ and radius $R = 0.15\;\text{m}$ is released from rest at height $H$ (measured from the bottom of a circular loop) on an inclined track. At the bottom of the incline the track forms a vertical circular loop of radius $r = 0.8\;\text{m}$. The disk rolls without slipping throughout.

Find the minimum height $H$ for the disk to maintain contact with the track at the top of the loop. \qmarks{5}
\end{parts}

\newpage

% ---- QUESTION C2: PROPERTIES OF MATTER (15 marks) ----
\question{15}

\noindent\textbf{Advanced Entropy, Lennard-Jones Potential, and Maxwell--Boltzmann Distribution}

\begin{parts}
\item Two identical blocks of the same material (mass $m$, specific heat capacity $c$) are initially at temperatures $T_1$ and $T_2$ ($T_1 > T_2$). A series of Carnot engines operate between them, each extracting an infinitesimal amount of work, until the blocks reach the same final temperature $T_f$.
\begin{parts}
\item Show that $T_f = \sqrt{T_1\,T_2}$. \qmarks{2}
\item Show that the total entropy change of the two blocks is zero. \qmarks{2}
\end{parts}

\item If instead the two blocks are simply placed in direct thermal contact (an irreversible process):
\begin{parts}
\item Show that the final temperature is $T_f' = \frac{T_1 + T_2}{2}$. \qmarks{1}
\item Calculate the total entropy change $\Delta S_\text{total}$ and show that it is positive. \qmarks{2}
\item Hence show that the maximum extractable work from the Carnot process in part~(a) was $W_\text{max} = mc\!\left(T_1 + T_2 - 2\sqrt{T_1 T_2}\right)$. \qmarks{1}
\end{parts}

\item The Lennard-Jones potential is $V(r) = 4\varepsilon\!\left[\left(\dfrac{a}{r}\right)^{12} - \left(\dfrac{a}{r}\right)^{6}\right]$.
\begin{parts}
\item Show that the equilibrium separation is $r_0 = 2^{1/6}\,a$ and that the depth of the potential well is $\varepsilon$. \qmarks{2}
\item By expanding $V(r)$ about $r_0$ to second order, show that the effective spring constant is $k = 72\,\varepsilon / (2^{1/3}\,a^2)$. Hence estimate the frequency of small oscillations of a single argon atom about its equilibrium position, using $\varepsilon = 1.67\times 10^{-21}\;\text{J}$, $a = 3.40\times 10^{-10}\;\text{m}$, and $m_\text{Ar} = 6.63\times 10^{-26}\;\text{kg}$. \qmarks{2}
\end{parts}

\item The most probable speed in the Maxwell--Boltzmann distribution is $c_\text{mp} = \sqrt{2k_BT/m}$. At what temperature does the most probable speed of helium atoms ($m_\text{He} = 6.64\times 10^{-27}\;\text{kg}$) equal the escape velocity from the Earth's surface ($v_e = 11.2\;\text{km/s}$)? Comment on whether helium can escape from the Earth's atmosphere at room temperature. \qmarks{3}
\end{parts}

\newpage

% ---- QUESTION C3: OSCILLATIONS AND WAVES (15 marks) ----
\question{15}

\noindent\textbf{Advanced Boundary Conditions, Beaded String, and Wave Packets}

\begin{parts}
\item A string of length $L$ is fixed at $x = 0$ and attached to a massless ring that slides freely (without friction) on a vertical rod at $x = L$.
\begin{parts}
\item What is the boundary condition on the displacement $y(x,t)$ at $x = L$? \qmarks{1}
\item Show that the allowed wavelengths are $\lambda_n = \dfrac{4L}{2n - 1}$ for $n = 1, 2, 3, \ldots$ \qmarks{2}
\item If $L = 2\;\text{m}$ and $v = 100\;\text{m/s}$, find the first three resonant frequencies. \qmarks{2}
\end{parts}

\item Two identical strings (same linear mass density $\mu$, same tension $T$) of length $L$ are connected at $x = 0$ by a bead of mass $m$. The left string is fixed at $x = -L$ and the right string is fixed at $x = +L$.
\begin{parts}
\item Write down the boundary conditions at $x = \pm L$ and the matching conditions at $x = 0$ (continuity of displacement and the force balance including the bead's inertia). \qmarks{3}
\item For the \emph{symmetric} normal modes, the bead sits at a node (i.e.\ $y(0,t) = 0$). Show that in this case the resonant frequencies are $f_n = nv/(2L)$ for $n = 1, 2, 3, \ldots$ \qmarks{2}
\end{parts}

\item A wave packet is formed from two harmonic waves:
\[
y(x,t) = A\cos(k_1 x - \omega_1 t) + A\cos(k_2 x - \omega_2 t),
\]
where $k_1 = 10\;\text{rad/m}$, $k_2 = 12\;\text{rad/m}$, and the dispersion relation is $\omega = \alpha k^2$ with $\alpha = 0.5\;\text{m}^2/\text{s}$.
\begin{parts}
\item Calculate $\omega_1$ and $\omega_2$. \qmarks{1}
\item Find the phase velocity $v_p$ and group velocity $v_g$ at the average wavenumber $\bar{k} = (k_1 + k_2)/2$. \qmarks{2}
\item Rewrite $y(x,t)$ as a product of a slowly-varying envelope and a rapidly-varying carrier. Identify which factor moves at the group velocity and which at the phase velocity. \qmarks{2}
\end{parts}
\end{parts}

\newpage

% ---- QUESTION C4: QUANTUM PHYSICS (15 marks) ----
\question{15}

\noindent\textbf{Three-State System, Fermion States, and Pion Decay}

\begin{parts}
\item Consider a three-state quantum system with orthonormal basis $\{\ket{a},\,\ket{b},\,\ket{c}\}$. The Hamiltonian is
\[
H = E_0\bigl(\ketbra{a}{a} + \ketbra{b}{b} + \ketbra{c}{c}\bigr) + V\bigl(\ketbra{a}{b} + \ketbra{b}{a}\bigr),
\]
where $E_0$ and $V$ are real constants.
\begin{parts}
\item Write $H$ as a $3\times 3$ matrix in the $\{\ket{a},\,\ket{b},\,\ket{c}\}$ basis. \qmarks{1}
\item Find all three eigenvalues and the corresponding normalised eigenstates. \qmarks{3}
\end{parts}

\item At $t = 0$ the system is prepared in the state $\ket{\psi(0)} = \ket{a}$.
\begin{parts}
\item Find $\ket{\psi(t)}$. \qmarks{2}
\item Calculate the probability $P(a,t) = |\braket{a}{\psi(t)}|^2$ of finding the system in state $\ket{a}$ at time $t$. \qmarks{2}
\end{parts}

\item Two identical spin-$\frac{1}{2}$ fermions occupy a system with two spatial orbitals $\ket{\alpha}$ and $\ket{\beta}$. Write down \emph{all} possible two-particle states that are antisymmetric under exchange of the two fermions (combining spatial and spin degrees of freedom). State the total number of such states. \qmarks{3}

\item Consider the decay $\pi^0 \to \gamma + \gamma$ (neutral pion decaying into two photons).
\begin{parts}
\item Verify that charge, baryon number, and lepton number are all conserved in this decay. \qmarks{2}
\item The $\pi^0$ has spin~0. Using conservation of angular momentum, discuss the constraint on the helicities (circular polarisation states) of the two emitted photons. Which force mediates this decay? \qmarks{2}
\end{parts}
\end{parts}

\newpage

% ============================================================
% FORMULA SHEET
% ============================================================
% EQUATION SHEET — PHYS10012 Core Physics I
% This file is \input{} at the end of each exam paper

\newpage
\section*{Formula Sheet}
\noindent\rule{\textwidth}{0.4pt}

\subsection*{Mathematical Formulae}

\begin{align*}
&\text{Binomial expansion:} & (1+x)^n &\approx 1 + nx + \frac{n(n-1)}{2!}x^2 + \cdots \quad (|x| \ll 1) \\[4pt]
&\text{Euler's formula:} & e^{i\theta} &= \cos\theta + i\sin\theta \\[4pt]
&\text{Trig identities:} & \sin A + \sin B &= 2\cos\!\left(\tfrac{A-B}{2}\right)\sin\!\left(\tfrac{A+B}{2}\right) \\
& & \cos A + \cos B &= 2\cos\!\left(\tfrac{A-B}{2}\right)\cos\!\left(\tfrac{A+B}{2}\right)
\end{align*}

\subsection*{Mechanics}

\begin{align*}
&\text{SUVAT:} & v &= u + at, \quad s = ut + \tfrac{1}{2}at^2, \quad v^2 = u^2 + 2as, \quad s = \tfrac{1}{2}(u+v)t \\[4pt]
&\text{Dot product:} & \mathbf{a}\cdot\mathbf{b} &= |\mathbf{a}||\mathbf{b}|\cos\theta = a_x b_x + a_y b_y + a_z b_z \\[4pt]
&\text{Cross product:} & \mathbf{a}\times\mathbf{b} &= |\mathbf{a}||\mathbf{b}|\sin\theta\;\hat{\mathbf{n}} = \begin{vmatrix}\hat{\mathbf{i}}&\hat{\mathbf{j}}&\hat{\mathbf{k}}\\a_x&a_y&a_z\\b_x&b_y&b_z\end{vmatrix} \\[4pt]
&\text{Rocket equation:} & \Delta v &= u \ln\!\left(\frac{m_0}{m_f}\right) \\[4pt]
&\text{Gravitational PE:} & U &= -\frac{GMm}{r} \\[4pt]
&\text{Escape velocity:} & v_e &= \sqrt{\frac{2GM}{R}} \\[4pt]
&\text{Force from PE:} & F &= -\frac{dU}{dx} \\[4pt]
&\text{Elastic collision (1D):} & v_1' &= \frac{m_1 - m_2}{m_1 + m_2}v_1 + \frac{2m_2}{m_1+m_2}v_2 \\[4pt]
&\text{Centre of mass:} & \mathbf{R}_\text{cm} &= \frac{\sum m_i\mathbf{r}_i}{\sum m_i} \\[4pt]
&\text{Reduced mass:} & \mu &= \frac{m_1 m_2}{m_1 + m_2} \\[4pt]
&\text{Centripetal accel:} & a_c &= \frac{v^2}{r} = \omega^2 r \\[4pt]
&\text{Parallel axis thm:} & I &= I_\text{cm} + Md^2 \\[4pt]
&\text{Atwood machine:} & a &= \frac{(m_1-m_2)g}{m_1+m_2}, \quad T = \frac{2m_1 m_2 g}{m_1+m_2}
\end{align*}

\noindent\textbf{Moments of inertia} (about axis through centre of mass):
\begin{center}
\begin{tabular}{ll}
Thin ring (radius $R$): & $I = MR^2$ \\
Solid disk (radius $R$): & $I = \tfrac{1}{2}MR^2$ \\
Solid sphere (radius $R$): & $I = \tfrac{2}{5}MR^2$ \\
Hollow sphere (radius $R$): & $I = \tfrac{2}{3}MR^2$ \\
Thin rod (length $L$, about centre): & $I = \tfrac{1}{12}ML^2$ \\
Thin rod (length $L$, about end): & $I = \tfrac{1}{3}ML^2$ \\
\end{tabular}
\end{center}

\subsection*{Properties of Matter}

\begin{align*}
&\text{Ideal gas:} & PV &= nRT \\[4pt]
&\text{Heat capacities:} & C_P - C_V &= R \quad\text{(per mole)} \\[4pt]
&\text{Adiabatic:} & PV^\gamma &= \text{const}, \quad TV^{\gamma-1} = \text{const} \\[4pt]
&\text{Isothermal work:} & W &= nRT\ln\!\left(\frac{V_2}{V_1}\right) \\[4pt]
&\text{Carnot efficiency:} & \eta &= 1 - \frac{T_c}{T_h} \\[4pt]
&\text{Entropy:} & dS &= \frac{dQ_\text{rev}}{T}, \quad \Delta S = mc\ln\!\left(\frac{T_2}{T_1}\right), \quad \Delta S = nR\ln\!\left(\frac{V_2}{V_1}\right) \\[4pt]
&\text{Lennard-Jones:} & V(r) &= 4\varepsilon\left[\left(\frac{a}{r}\right)^{12} - \left(\frac{a}{r}\right)^{6}\right] \\[4pt]
&\text{Mean KE:} & \langle KE \rangle &= \tfrac{3}{2}k_B T \\[4pt]
&\text{MB speeds:} & c_\text{mp} &= \sqrt{\frac{2k_BT}{m}}, \quad c_\text{avg} = \sqrt{\frac{8k_BT}{\pi m}}, \quad c_\text{rms} = \sqrt{\frac{3k_BT}{m}}
\end{align*}

\subsection*{Oscillations and Waves}

\begin{align*}
&\text{SHM:} & x(t) &= A\cos(\omega_0 t + \phi), \quad \omega_0 = \sqrt{k/m} \\[4pt]
&\text{Damped SHM:} & x(t) &= Ae^{-\gamma t/2}\cos(\omega_d t + \phi), \quad \omega_d = \sqrt{\omega_0^2 - \gamma^2/4} \\[4pt]
&\text{Quality factor:} & Q &= \omega_0/\gamma \\[4pt]
&\text{Forced amplitude:} & A(\omega) &= \frac{F_0/m}{\sqrt{(\omega_0^2 - \omega^2)^2 + \gamma^2\omega^2}} \\[4pt]
&\text{Phase lag:} & \tan\delta &= \frac{\gamma\omega}{\omega_0^2 - \omega^2} \\[4pt]
&\text{Pendulums:} & T_\text{simple} &= 2\pi\sqrt{L/g}, \quad T_\text{physical} = 2\pi\sqrt{I/(Mgh)} \\[4pt]
&\text{Wave equation:} & \frac{\partial^2 y}{\partial x^2} &= \frac{1}{v^2}\frac{\partial^2 y}{\partial t^2} \\[4pt]
&\text{String speed:} & v &= \sqrt{T/\mu} \\[4pt]
&\text{Sound in gas:} & v &= \sqrt{\gamma P/\rho} \\[4pt]
&\text{Harmonic wave:} & y &= A\sin(kx - \omega t) \\[4pt]
&\text{String impedance:} & Z &= \mu v = \sqrt{\mu T} \\[4pt]
&\text{Reflection:} & r &= \frac{Z_1 - Z_2}{Z_1 + Z_2}, \quad t = \frac{2Z_1}{Z_1+Z_2} \\[4pt]
&\text{Power (string):} & \langle P\rangle &= \tfrac{1}{2}\mu\omega^2 A^2 v \\[4pt]
&\text{Decibels:} & \beta &= 10\log_{10}(I/I_0), \quad I_0 = 10^{-12}\;\text{W/m}^2 \\[4pt]
&\text{Doppler:} & f' &= f\frac{v \pm v_o}{v \mp v_s} \\[4pt]
&\text{Group velocity:} & v_g &= \frac{d\omega}{dk} \\[4pt]
&\text{String harmonics:} & f_n &= \frac{nv}{2L} \quad (n=1,2,3,\ldots) \\[4pt]
&\text{Open-closed pipe:} & f_n &= \frac{nv}{4L} \quad (n=1,3,5,\ldots)
\end{align*}

\subsection*{Quantum Physics}

\begin{align*}
&\text{Schr\"{o}dinger eq:} & i\hbar\frac{d}{dt}\ket{\psi(t)} &= H\ket{\psi(t)} \\[4pt]
&\text{Time evolution:} & \ket{\psi(t)} &= \sum_n c_n e^{-iE_n t/\hbar}\ket{E_n}, \quad U(t) = e^{-iHt/\hbar} \\[4pt]
&\text{Spin operators:} & S_z &= \frac{\hbar}{2}\bigl(\ketbra{\!\uparrow}{\uparrow\!} - \ketbra{\!\downarrow}{\downarrow\!}\bigr) \\
& & S_x &= \frac{\hbar}{2}\bigl(\ketbra{\!\uparrow}{\downarrow\!} + \ketbra{\!\downarrow}{\uparrow\!}\bigr) \\
& & S_y &= \frac{\hbar}{2}\bigl(-i\ketbra{\!\uparrow}{\downarrow\!} + i\ketbra{\!\downarrow}{\uparrow\!}\bigr) \\[4pt]
&\text{Spin eigenstates:} & \ket{\!\uparrow_x} &= \tfrac{1}{\sqrt{2}}\bigl(\ket{\!\uparrow_z} + \ket{\!\downarrow_z}\bigr), \quad \ket{\!\downarrow_x} = \tfrac{1}{\sqrt{2}}\bigl(\ket{\!\uparrow_z} - \ket{\!\downarrow_z}\bigr) \\
& & \ket{\!\uparrow_y} &= \tfrac{1}{\sqrt{2}}\bigl(\ket{\!\uparrow_z} + i\ket{\!\downarrow_z}\bigr), \quad \ket{\!\downarrow_y} = \tfrac{1}{\sqrt{2}}\bigl(\ket{\!\uparrow_z} - i\ket{\!\downarrow_z}\bigr) \\[4pt]
&\text{Pauli matrices:} & \sigma_x &= \begin{pmatrix}0&1\\1&0\end{pmatrix}, \quad \sigma_y = \begin{pmatrix}0&-i\\i&0\end{pmatrix}, \quad \sigma_z = \begin{pmatrix}1&0\\0&-1\end{pmatrix}
\end{align*}

\subsection*{Physical Constants}

\begin{center}
\begin{tabular}{lll}
Speed of light & $c$ & $3.00 \times 10^{8}$ m/s \\
Planck's constant & $h$ & $6.626 \times 10^{-34}$ J$\cdot$s \\
Reduced Planck's constant & $\hbar$ & $1.055 \times 10^{-34}$ J$\cdot$s \\
Boltzmann constant & $k_B$ & $1.381 \times 10^{-23}$ J/K \\
Gas constant & $R$ & $8.314$ J/(mol$\cdot$K) \\
Avogadro's number & $N_A$ & $6.022 \times 10^{23}$ mol$^{-1}$ \\
Gravitational constant & $G$ & $6.674 \times 10^{-11}$ N$\cdot$m$^2$/kg$^2$ \\
Acceleration due to gravity & $g$ & $9.81$ m/s$^2$ \\
Electron mass & $m_e$ & $9.109 \times 10^{-31}$ kg \\
Proton mass & $m_p$ & $1.673 \times 10^{-27}$ kg \\
Elementary charge & $e$ & $1.602 \times 10^{-19}$ C \\
Mass of Earth & $M_\oplus$ & $5.97 \times 10^{24}$ kg \\
Radius of Earth & $R_\oplus$ & $6.37 \times 10^{6}$ m \\
\end{tabular}
\end{center}

\noindent\rule{\textwidth}{0.4pt}


\end{document}
